\section{Discussion}\label{sec:discussion}

\paragraph{What a boundary map could show}
The public argument about Israel-related speech is usually conducted as a
single question: is this kind of criticism antisemitic? The framework
replaces it with many narrower questions, one for each statement in a
pre-registered item bank, each with an answer that has an interval and
stated assumptions. For each Israel-related statement, a boundary map would
show how much of the common variance in responses to it is shared, through
the general factor, with responses to classic stereotypes
(\cref{def:boundary}), and whether that share lies clearly above, clearly
below or indeterminately around a threshold fixed in advance
(\cref{fig:boundary,prop:id-boundary-test}). It would show this only where the data identify it
(\cref{thm:ident}), and it would say so where they do not. The
correlated-factors model adds a coarser summary, the latent correlation of
the two attitude domains. It is the quantity closest to the scale-level
associations reported in earlier studies (\cref{sec:bg-prior}), and it
cannot by itself separate a common cause from correlated but distinct
attitudes (\cref{prop:model-sl}); the item-level map can, where the items
within a domain differ in their general-factor share.

Every outcome of such a map is informative to every side of the debate.
Someone who holds that a category of Israel-related speech typically
expresses hostility toward Jews learns whether, in the surveyed population,
statements in that category covary with classic stereotypes as that view
implies. Someone who holds that the category typically expresses a distinct
political attitude learns the same thing about the opposite claim. Both learn
where the evidence is indeterminate at the sample size used, which is itself
a finding. Under partial invariance the map can also be compared across
political, generational and national groups (\cref{thm:dif-bias,prop:anchor}).
These comparisons can be anchored in the stereotype items alone
(\cref{rem:id-anchor-S}), so no one need assume that any Israel-related
statement means the same thing to the groups being compared, an assumption
that would itself be contested.
An item could sit inside the boundary set in one group and outside it in
another. A finding of that kind would bear directly on whether one rule can
be applied to speakers from every group, a question that the definitions
answer differently through the weight they give to context
(\cref{sec:bg-definitions}). The list-experiment module would show whether
direct answers understate or overstate endorsement, item by item
(\cref{thm:joint-sdb}). Under a uniform threshold shift, social desirability
biases prevalences but leaves the boundary estimands unchanged
(\cref{prop:list-irt}); it affects latent-mean comparisons only if the shift
differs between groups, where it acts as DIF. The map is therefore robust to
the simplest form of strategic answering, though not to a shift that acts
differently on the general and the specific parts of an item's response. Multilevel regression and
poststratification would carry the latent trait to regions and states with
honest intervals (\cref{eq:mrp}), and \cref{prop:index} explains why a
stereotype-count index cannot do the same job without invariance
diagnostics.

\paragraph{What it could not show}
A boundary map is not a verdict. It does not say which speech ought to be
called antisemitic, and it is compatible with more than one answer. A
definition could reasonably treat a statement as antisemitic even if its
$\mathrm{ECV}_j$ is low, for example because of what the statement does to
the people it is about, or because definitions classify speech in context.
A definition could equally decline to treat a statement as antisemitic even
if its $\mathrm{ECV}_j$ is high, because covariance with stereotypes across
a population does not establish what any particular speaker meant.

Nor is the map a diagnosis of anyone. $\mathrm{ECV}_j$ is a ratio of
variances: it says how responses to an item move with the general factor
across the population, not how many of the item's endorsers stand high on
that factor (\cref{sec:model-meaning}). An item can have a high share and
still be endorsed by many respondents who stand low on the general factor,
if its threshold is low. Nor is the map causal. A general factor summarises shared variance, and shared variance can have
several sources: a common attitude, a common information environment, or a
common response style. The design separates some of these, through
knowledge measures, reverse-keyed items and comparison-state items
(\cref{sec:design-items}), but not all. Finally, the general factor means
only what its items make it mean. We read it as hostility toward Jews as
Jews because the stereotype items that define it are not themselves in
dispute between the definitions. If those items also share a broader
disposition, such as a general readiness to endorse conspiratorial claims,
the general factor absorbs it, and the boundary estimands inherit that
ambiguity. This is why the stereotype battery, and not only the
Israel-related items, is fixed before fieldwork and reviewed by advisers with
opposing views (\cref{sec:ethics-advisory}). Whether the general factor
should be defined as what both domains share, as in our primary model, or
anchored in the stereotype domain alone, as in the bifactor-$(S-1)$ variant
reported as a sensitivity analysis (\cref{rem:id-s1}), is itself part of the
substantive question.

\paragraph{Limitations}
We stress five.

\emph{Survey mode.} The plan uses a self-administered online panel
(\cref{sec:design-items}). Mode affects what respondents are willing to
report \citep{kreuter2008,tourangeau2007}, panels differ in coverage, and the
assumptions of the list experiment, no design effect and no liars
(\cref{asm:le}), can be tested only in part (\cref{cor:le-tests}). Estimates
from one mode need not transfer to another, and the social-desirability
correction is specific to the mode in which it is estimated.

\emph{Dependence on item wording.} $\mathrm{ECV}_j$ is a property of an item,
a particular wording answered by a particular population, and not of a
category of speech. Two wordings of the same idea can load differently, and a
statement that is double-barrelled, ambiguous or unfamiliar can land
anywhere. The map generalises from items to categories only as far as the
bank samples the wordings of each category, which is why the plan requires
several items per disputed category and records how each definition treats
each item (\cref{sec:design-items}). A claim about a category resting on one
item would be a claim about that item.

\emph{Translation across countries and cultures.} A statement about Israel
does not mean the same thing in every country, language or political
culture. Translation can create differential item functioning, and scalar
invariance is often hard to establish across countries
\citep{davidov2014,steenkamp1998}. A cross-national boundary map therefore
needs a translation protocol, invariance testing of every item in every
country, and, where invariance fails, anchoring vignettes \citep{king2004} or
alignment \citep{asparouhov2014}. Without these, a difference in boundary
sets between countries could reflect translation rather than attitudes.

\emph{Latent-variable assumptions.} Every result rests on a model: a graded
response link, local independence given the factors, normal latent
distributions, orthogonal factors in the bifactor model, and an assignment of
each item to exactly one specific factor. If an Israel-related item also
loads on the stereotype factor, a bifactor model without that cross-loading
will tend to credit the shared variance to the general factor. Our
identification results state what is identified under the model, not that
the model is true. The decision rules compare the bifactor model with the
correlated-factors and unidimensional models and report a boundary set only
when the bifactor model passes pre-registered fit criteria
(\cref{sec:design-decisions}). The simulation study tests the estimators
under the departures it simulates (\cref{sec:simulation}), not under every
departure that real data might present.

\emph{Population association is not individual intent.} The estimands
describe how attitudes covary across a population at the time of fieldwork.
They say nothing about the intentions of any respondent, speaker or
organisation, and they may change as events change what statements about
Israel mean to the people answering them. Every output will carry these
limits (\cref{sec:ethics-individual}).

\paragraph{Extensions}
Three extensions seem most valuable. First, \emph{longitudinal invariance}.
Attitudes about Israel may move with events, and so may the meaning of the
items. Tracking a boundary map over time requires invariance across waves as
well as groups, with anchor items that are themselves stable over time; an
item that drifts across the boundary could be changing in meaning, not in
the attitudes behind it. The event calendars of the second paper in this
series \citep{kornsteinlab2026contagion} offer a principled way to place
waves before and after salient external events. Second, \emph{conjoint
designs}. The item-level estimand is observational: it describes how
responses to fixed wordings covary. A conjoint design that randomises
attributes of a statement, such as whether it refers to Israel or to another
state, whether it invokes a classic trope, and how it is framed, would
estimate the average causal effect of each attribute on endorsement, the
average marginal component effect of \citet{hainmueller2014}. That would separate the content of a statement from
its wording, and so address the dependence on wording described above.
Third, \emph{linkage to text}. The lab's other projects propose to measure
antisemitism in online discourse, to evaluate detection algorithms and to
trace narratives \citep{kornsteinlab2026measurement,kornsteinlab2026detection,kornsteinlab2026narratives}.
A boundary map could inform them: the definitional tags and estimated
$\mathrm{ECV}_j$ of each item give annotators and classifiers a graded,
population-based description of Israel-related content in place of a single
binary label. In the other direction, narratives observed in text can
suggest wordings for the item bank. Any linkage of survey respondents to
their own online text would require separate consent and its own ethics
review (\cref{sec:ethics}).

\paragraph{Measurement as common ground}
The debate over where criticism of Israel ends and antisemitism begins will
not be settled by a survey. It turns on values, law and the experience of
people on every side, and those are not measured by covariance. But the
debate also contains empirical claims, and those can be tested in a way that
does not depend on which side one is on. The threshold, the items, the
decision rules and the code are fixed before the data are seen; the same
rule is applied to every item; and every estimate is published whatever its
direction. A result produced that way is not a weapon for either side. It is
a shared reference point that people who disagree about definitions can
inspect, criticise and replicate. That is what measurement can offer a
contested question: not an answer to it, but common ground on which to argue
about it.
